\documentclass[conference]{IEEEtran}

\usepackage[utf8]{inputenc}
\usepackage[T1]{fontenc}
\usepackage[spanish, es-nodecimaldot]{babel}

\usepackage{amsmath}
\usepackage{amssymb}
\usepackage{array}
\usepackage{booktabs}
\usepackage{multirow}

\usepackage{graphicx}

\usepackage{url}
\usepackage{cite}

\spanishdecimal{.}

\begin{document}

\title{Análisis de Variables Climáticas y Producción Agrícola en el Perú mediante un Enfoque Multi-distrito (2020--2025)}

\author{
\IEEEauthorblockN{Johao Mendoza, Joyssie Rivas, Thiago Ormeño}
\IEEEauthorblockA{\textit{Departamento de Ingeniería} \\
\textit{Universidad del Pacífico} \\
Lima, Perú \\
\{jr.mendozaf@alum.up.edu.pe, ja.rivasr@alum.up.edu.pe, tc.ormenof@alum.up.edu.pe\}}
}

\maketitle

\begin{abstract}
Este estudio analiza la relación entre variables climáticas y producción agrícola en seis regiones del Perú (Ica, Junín, La Libertad, Piura, Puno y San Martín) durante 2020--2025, integrando producción mensual del MIDAGRI con variables climáticas de NASA POWER. Debido a la heterogeneidad ecológica intra-regional, se implementó una estrategia multi-distrito basada en 14 unidades representativas por piso ecológico y asociadas a cultivos mediante criterio agronómico. La producción mensual fue recuperada desde los acumulados oficiales del MIDAGRI mediante diferenciación temporal y manejo explícito de datos faltantes. Los resultados muestran que los determinantes climáticos varían según el piso ecológico: la radiación domina en costa, la humedad de suelo y temperatura máxima en sierra y altiplano, y la humedad relativa en selva. La sequía altoandina de 2022 en Puno redujo fuertemente la producción de quinua ($-64\%$) y oca ($-62\%$), mientras que el Fenómeno del Niño 2023--2024 produjo anomalías regionales diferenciadas, destacando el incremento de precipitación en Piura (+121\%). Los resultados validan la utilidad del enfoque multi-distrito para capturar heterogeneidad agroclimática que permanece oculta bajo agregados regionales convencionales.
\end{abstract}

\begin{IEEEkeywords}
Agricultura, cambio climático, NASA POWER, MIDAGRI, análisis agroclimático, piso ecológico, ENSO, CRISP-DM
\end{IEEEkeywords}

\section{Introducción}
La agricultura representa el 4.8\% del PBI peruano en 2025 \cite{bcrp2025} y emplea aproximadamente 4.3 millones de personas (24\% de la PEA) \cite{inei2024}. Las agroexportaciones crecieron de USD 7,791 millones en 2020 a USD 15,066 millones en 2025 \cite{adex2026}. La producción exhibe alta variabilidad por eventos climáticos (sequías, heladas, El Niño) cuya relación con los rendimientos por cultivo no ha sido analizada sistemáticamente con técnicas de Data Mining a granularidad mensual.

La complejidad metodológica del análisis agroclimático en el Perú es estructural: las regiones son ecológicamente heterogéneas (una sola región puede combinar costa, sierra y selva), y los datos públicos de producción del MIDAGRI se reportan a nivel regional acumulado, no por piso ecológico ni a granularidad mensual directa. Trabajos previos enfrentan estas limitaciones usando agregados anuales \cite{senamhi2009, unap2016}, asumiendo un único punto climático por región, o restringiéndose a un solo cultivo. La presente investigación combina tres aportes: (i) recuperación de producción mensual desde los acumulados publicados, (ii) muestreo climático multi-distrito por piso ecológico, y (iii) integración cultivo-piso con criterio agronómico documentado.

El objetivo general es caracterizar la relación entre variables climáticas y producción agrícola en seis regiones del Perú bajo la metodología CRISP-DM \cite{inia2019}. Los objetivos específicos son: 
\begin{enumerate}
    \item Construir un dataset integrado clima-producción a granularidad cultivo-piso.
    \item Identificar los determinantes climáticos por región y piso ecológico.
    \item Detectar patrones de vulnerabilidad y estacionalidad por cultivos agrícolas.
    \item Validar la utilidad metodológica del enfoque multi-distrito mediante el caso de la sequía altoandina 2022 en Puno y el Fenómeno del Niño 2023-2024 en regiones costeras.
\end{enumerate}

\section{Estado del Arte}
Diversos estudios han analizado la relación entre variabilidad climática y producción agrícola en el Perú, especialmente frente a eventos ENSO, sequías y cambios de temperatura y precipitación. Restrepo et al. \cite{restrepo2017} evidenciaron que eventos extremos de El Niño alteran significativamente patrones ambientales e hidrológicos en Sudamérica, afectando sistemas agrícolas vulnerables. Asimismo, estudios nacionales destacan diferencias agroclimáticas entre costa, sierra y selva debido a la heterogeneidad ecológica del territorio peruano.

En paralelo, el uso de metodologías de Data Mining y CRISP-DM en agricultura ha crecido en los últimos años. Chapapría y Viladrich \cite{chapapria2024} desarrollaron un modelo agroclimático basado en CRISP-DM, mientras que Canchano \cite{canchano2024} aplicó técnicas predictivas para productos agrícolas. Sin embargo, estos enfoques se centran principalmente en tareas predictivas y presentan limitaciones en integración espacial y temporal de datos.

Por otro lado, trabajos de zonificación agroecológica desarrollados por MIDAGRI \cite{midagri2025} y la Universidad Nacional del Altiplano \cite{una2023} resaltan la importancia de considerar pisos ecológicos en estudios agrícolas. No obstante, gran parte de la literatura revisada emplea agregados anuales, análisis de un solo cultivo o un único punto climático por región, limitando la representación de microclimas y dinámicas intra-regionales.

En consecuencia, persiste un vacío metodológico en la integración de producción agrícola mensual, variables climáticas multi-distrito y caracterización agroecológica para múltiples regiones peruanas bajo un enfoque de minería de datos.


\section{Metodología}

\subsection{Fuentes de datos y variables}

Se adoptó la metodología CRISP-DM \cite{chapman2000} en las fases de comprensión del negocio, comprensión y preparación de datos, y análisis exploratorio. Se integraron dos fuentes públicas: (i) la hoja \texttt{c-18} de los reportes mensuales del MIDAGRI \cite{midagri2025}, que contiene producción agrícola acumulada por región, mes y cultivo para el periodo 2020--2025; y (ii) variables agroclimáticas de NASA POWER \cite{nasa2025}. 

Las variables climáticas fueron seleccionadas por su relevancia sobre fenología, disponibilidad hídrica y rendimiento agrícola \cite{viladrich}. La Tabla~\ref{tab:variables} resume las variables utilizadas.

\begin{table}[h]
\centering
\caption{Variables climáticas seleccionadas}
\label{tab:variables}
\small
\begin{tabular}{l l}
\toprule
\textbf{Variable} & \textbf{Descripción} \\
\midrule
T2M & Temperatura media \\
T2M\_MAX & Temperatura máxima \\
T2M\_MIN & Temperatura mínima \\
PRECTOTCORR & Precipitación corregida \\
RH2M & Humedad relativa \\
ALLSKY\_SFC & Radiación solar \\
\bottomrule
\end{tabular}
\end{table}

\subsection{Recuperación y calidad de los datos}

La hoja \texttt{c-18} reporta producción acumulada anual y no producción mensual directa. La producción mensual se recuperó mediante diferenciación entre meses consecutivos:

\begin{equation}
Prod_{m} = Acum_{m} - Acum_{m-1}
\end{equation}

Tres meses no publicados por MIDAGRI (dic-2020, may-2021 y mar-2022) fueron incorporados como valores faltantes explícitos antes del cálculo de diferencias para evitar sobreestimaciones posteriores. Diciembre 2020 se recuperó a partir del total anual publicado en el archivo diciembre 2021.

La Tabla~\ref{tab:calidad} resume los principales tratamientos de calidad aplicados.

\begin{table}[h]
\centering
\caption{Tratamientos de calidad de datos}
\label{tab:calidad}
\small
\begin{tabular}{p{2cm} p{4.3cm}}
\toprule
\textbf{Caso} & \textbf{Tratamiento} \\
\midrule
Meses faltantes MIDAGRI & NaN explícito y recuperación Dic-2020 \\
Valores \texttt{-999} NASA & Tratados como NaN y estimados con continuidad temporal local. \\
Negativos en \texttt{diff()} & Conversión a NaN \\
Ceros estructurales & Conservados \\
Eventos extremos & Conservados como señal \\
\bottomrule
\end{tabular}
\end{table}

No se aplicó remoción de outliers mediante IQR o Z-score debido a la alta presencia de ceros estructurales y a que eventos extremos como la sequía 2022 en Puno \cite{rcr2023} y El Niño Costero 2023--2024 \cite{senamhi2024} constituyen la señal de interés del estudio.

\subsection{Estrategia multi-distrito y construcción de datasets}

Debido a la heterogeneidad ecológica de las regiones peruanas, se seleccionaron 14 distritos representativos distribuidos por piso ecológico para evitar sesgos espaciales derivados del uso de un único punto climático regional. La selección consideró especialización agrícola y representatividad agroecológica.

\begin{table}[h]
\centering
\caption{Distritos representativos por piso ecológico}
\label{tab:distritos}
\small
\begin{tabular}{l l l}
\toprule
\textbf{Región} & \textbf{Distrito} & \textbf{Piso ecológico} \\
\midrule
Ica & Chincha Alta & Costa \\
La Libertad & Virú & Costa \\
La Libertad & Huamachuco & Sierra \\
La Libertad & Cascas & Yunga \\
Piura & Tambogrande & Bosque seco \\
Piura & Sullana & Valle Chira \\
Piura & Canchaque & Sierra cafetalera \\
San Martín & Moyobamba & Selva Alto Mayo \\
San Martín & Tocache & Selva Huallaga \\
Junín & Perené & Selva alta \\
Junín & Río Tambo & Selva baja \\
Junín & El Tambo & Sierra Mantaro \\
Puno & Ilave & Altiplano lacustre \\
Puno & Ayaviri & Puna alta \\
\bottomrule
\end{tabular}
\end{table}

{\footnotesize Fuente: elaboración propia basada en \cite{redesarrolloica2021, esan2023, cepiura2023, dataperu2021, inei1996, midagripuno2020}.}

Finalmente, se construyeron tres datasets: (i) un dataset regional agregado por región y distrito; (ii) un dataset clima-producción por cultivo; y (iii) un dataset filtrado mediante Pareto-80, conservando los cultivos que representan el 80\% de la producción regional total.


\section{Resultados}

\subsection{Concentración productiva y estacionalidad por piso}

La producción agregada 2020--2025 muestra una fuerte concentración en pocas unidades región-piso (Fig.~\ref{fig:prod_unidad}). La Libertad-Costa (Virú) lidera con 36.84 M tn, dominada por caña de azúcar agroindustrial. Le siguen las dos unidades altoandinas de Puno: puna alta (Ayaviri) con 20.36 M tn y altiplano lacustre (Ilave) con 10.19 M tn, ambas asociadas principalmente a avena forrajera y alfalfa. Ica-Costa (Chincha Alta), con 9.32 M tn, exhibe el perfil agroexportador más diversificado (uva, mandarina, espárrago y palta). El rango entre la unidad más productiva y la menos productiva (Cascas-Yunga, 0.36 M tn) abarca dos órdenes de magnitud, evidenciando la heterogeneidad estructural del agro peruano y justificando el análisis por piso ecológico.

\begin{figure}[h]
    \centering
    \includegraphics[width=0.46\textwidth]{ProduccionPorUnidad.png}
    \caption{Producción acumulada 2020--2025 por unidad región-piso, en millones de toneladas.}
    \label{fig:prod_unidad}
\end{figure}

El calendario mensual de cosechas (Fig.~\ref{fig:calendario}) revela contrastes marcados de estacionalidad. Las unidades altoandinas de Puno concentran su producción entre febrero y mayo, con mínimos entre julio y septiembre, reflejando el ciclo único de cultivos andinos. En contraste, La Libertad-Costa e Ica-Costa mantienen producción relativamente uniforme durante el año, característica de sistemas agroexportadores intensivos y cultivos perennes. Las unidades de selva presentan estacionalidad intermedia.

\begin{figure}[h]
    \centering
    \includegraphics[width=0.46\textwidth]{CalendarioCosechas.png}
    \caption{Calendario mensual de producción por unidad región-piso.}
    \label{fig:calendario}
\end{figure}

\subsection{Relaciones clima-cultivo por piso ecológico}

Las correlaciones clima-cultivo más fuertes (Fig.~\ref{fig:corr_clima}) muestran patrones coherentes con la fisiología y condiciones agroecológicas de cada piso. En la costa iqueña, la radiación solar domina la estacionalidad: mandarina ($r=-0.78$) y maíz amarillo duro ($r=-0.58$) concentran cosechas en periodos de menor radiación, mientras que uva ($r=+0.61$) y tomate ($r=+0.55$) presentan comportamiento opuesto. La uva además correlaciona positivamente con humedad relativa ($r=+0.62$) y temperatura mínima ($r=+0.55$).

En Junín, la humedad relativa es dominante en cultivos de selva alta como piña ($r=+0.76$) y alfalfa ($r=+0.70$). La alfalfa también depende de humedad de suelo (GWETROOT, $r=+0.67$) y presenta sensibilidad negativa frente a temperatura máxima ($r=-0.51$). Otros cultivos andinos muestran sensibilidad térmica similar: papa ($r=-0.52$) y avena forrajera ($r=-0.52$) frente a T2M\_MAX, y naranja frente a heladas (T2M\_MIN, $r=-0.57$). En Puno, la alfalfa replica el patrón forrajero asociado a humedad de suelo ($r=+0.66$).

Las correlaciones variaron en magnitud y signo entre pisos ecológicos para variables climáticas equivalentes, evidenciando que los agregados regionales convencionales pueden ocultar relaciones agroclimáticas relevantes.

\begin{figure}[h]
    \centering
    \includegraphics[width=0.46\textwidth]{CorrelacionesClimaCultivo.png}
    \caption{Correlaciones de Pearson entre variables climáticas y producción por cultivo dominante.}
    \label{fig:corr_clima}
\end{figure}

\subsection{Eventos climáticos extremos: sequía 2022 y El Niño 2023--2024}

El año 2022 registró en Puno la precipitación más baja del periodo analizado: 1.36 mm/día en altiplano lacustre y 1.34 mm/día en puna alta, frente a aproximadamente 2.0 mm/día en años normales. La producción agregada cayó en consecuencia: altiplano lacustre registró una reducción cercana al 50\% respecto a 2020 y puna alta una caída del 37\%.

Sin embargo, el impacto fue más severo en cultivos andinos tradicionales. La Fig.~\ref{fig:puno_relativo} muestra que quinua cayó 64\%, oca 62\%, cebada forrajera 55\% y papa 47\% respecto al año base 2020. En contraste, la avena forrajera mostró menor sensibilidad hídrica, amortiguando parcialmente la caída del agregado regional. Estos resultados validan la necesidad del análisis desagregado por cultivo frente a agregados regionales.

\begin{figure}[h]
    \centering
    \includegraphics[width=0.46\textwidth]{PunoRelativo2020.png}
    \caption{Producción anual de cultivos puneños relativa a 2020. La franja sombreada indica la sequía de 2022.}
    \label{fig:puno_relativo}
\end{figure}

De forma complementaria, el periodo 2023--2024 registró anomalías térmicas y de precipitación consistentes con el Fenómeno del Niño Costero. Todas las regiones incrementaron su temperatura media entre 4.8\% y 9.2\% respecto al periodo 2020--2022, siendo Puno el caso más extremo. Piura presentó la mayor anomalía de precipitación (+121\%), coherente con la intensificación de lluvias en la costa norte reportada por SENAMHI \cite{senamhi2024}. En contraste, San Martín redujo su precipitación (-28\%) y humedad relativa (-8.6\%), evidenciando respuestas climáticas regionales diferenciadas. Estos resultados validan la capacidad del enfoque multi-distrito para capturar heterogeneidad espacial y climática que permanecería oculta bajo agregados regionales convencionales.

\begin{figure}[h]
    \centering
    \includegraphics[width=0.48\textwidth]{FenomenoNino2024.png}
    \caption{Producción, temperatura y precipitación promedio durante el periodo 2023--2024 respecto al baseline 2020--2022. Piura presenta la mayor anomalía de precipitación, mientras que todas las regiones muestran incremento térmico.}
    \label{fig:nino2024}
\end{figure}


\section{Discusión y conclusiones}

Los resultados muestran que el agro peruano responde a dinámicas climáticas diferenciadas según el piso ecológico, más que según la unidad administrativa regional. La radiación domina la estacionalidad costera, mientras que la humedad de suelo y la temperatura máxima condicionan la productividad en sierra y altiplano; en selva, la humedad relativa adquiere mayor relevancia. Estas relaciones incluso cambian de signo entre pisos para variables equivalentes, comportamiento que permanecería oculto bajo promedios regionales convencionales.

Los eventos extremos analizados refuerzan esta heterogeneidad. La sequía altoandina de 2022 redujo la producción agregada de Puno entre 37\% y 50\%, pero el análisis desagregado mostró impactos mucho más severos en cultivos andinos tradicionales como quinua ($-64\%$), oca ($-62\%$) y papa ($-47\%$), parcialmente amortiguados por la mayor resiliencia de la avena forrajera. De manera complementaria, el Fenómeno del Niño 2023--2024 produjo respuestas regionales contrastantes: Piura registró anomalías extremas de precipitación (+121\%), mientras que San Martín redujo precipitación y humedad relativa pese al incremento térmico generalizado observado en todas las regiones. Estos resultados evidencian que los eventos climáticos no afectan homogéneamente al territorio peruano y validan la utilidad del enfoque multi-distrito para capturar dicha variabilidad espacial.

El estudio presenta cuatro limitaciones principales. Primero, la asociación cultivo--distrito se construyó con criterio agronómico documentado, pero no puede validarse directamente debido a que MIDAGRI publica la producción solo a nivel regional. Segundo, la resolución espacial de NASA POWER ($\sim$50 km) limita la captura de microclimas, especialmente en zonas andinas de topografía compleja \cite{ramirez2012}. Tercero, el análisis se realizó en toneladas físicas y no en valor económico, subponderando cultivos de alta rentabilidad relativa. Finalmente, los reportes del MIDAGRI contienen revisiones retroactivas sobre datos preliminares (``p/''), lo cual introduce incertidumbre residual en las series recientes.

Como trabajo futuro, los datasets construidos permiten desarrollar modelos predictivos por cultivo-piso usando Random Forest o \textit{gradient boosting}, incorporar variables macroclimáticas asociadas a ENSO y estimar impactos económicos mediante integración con precios FOB agroexportadores \cite{chlingaryan2018}. La principal conclusión del estudio es que la región administrativa no constituye la unidad analítica natural del agro peruano: la heterogeneidad intra-regional exige enfoques de muestreo y análisis que respeten la estructura ecológica del territorio. Esto sugiere además que las políticas agrarias frente al cambio climático deberían diseñarse por piso ecológico y no únicamente por delimitaciones regionales.


\begin{thebibliography}{00} 

\bibitem{bcrp2025} Banco Central de Reserva del Perú, ``Actividad económica: Enero 2025,'' Nota de Estudios 23-2025, Mar. 2025. [Online]. Available: \url{https://www.bcrp.gob.pe/docs/Publicaciones/Notas-Estudios/2025/nota-de-estudios-23-2025.pdf} 

\bibitem{adex2026} Asociación de Exportadores (ADEX), ``Adex Hoy: Agroexportaciones 2025,'' Ed. 37081, Mar. 2026. [Online]. Available: \url{https://www.adexperu.org.pe/img/adexhoy/adexServAsopdf_37081_11-03-2026_18-01-25.pdf} 

\bibitem{inei2024} Instituto Nacional de Estadística e Informática, ``Encuesta Nacional Agraria,'' 2024. [Online]. Available: \url{https://www.inei.gob.pe/media/MenuRecursivo/publicaciones_digitales/Est/Lib1829/} 

\bibitem{inia2019} H. Dorado, ``Uso de minería de datos en agricultura para hacer frente al cambio climático,'' Instituto Nacional de Innovación Agraria (INIA), 2019. [Online]. Available: \url{https://www.inia.gob.pe/wp-content/uploads/2019/10/EventCC_Hugo_Dorado.pdf} 

\bibitem{senamhi2009} Servicio Nacional de Meteorología e Hidrología del Perú, ``Impacto del cambio climático en cultivos anuales,'' 2009. [Online]. Available: \url{https://web2.senamhi.gob.pe/load/file/01401SENA-12.pdf} 

\bibitem{unap2016} Universidad Nacional del Altiplano, ``Estudio agroclimatológico,'' Repositorio UNAP, 2016. [Online]. Available: \url{https://repositorio.unap.edu.pe/bitstreams/728acd12-860d-4a89-8876-8a6f2945484c/download} 

\bibitem{restrepo2017} J. M. Restrepo \textit{et al.}, ``The impact of extreme El Niño events on modern sediment transport and implications for the past,'' \textit{Scientific Reports}, vol. 7, no. 1, p. 11311, 2017, doi: 10.1038/s41598-017-12220-x. 

\bibitem{chapapria2024} L. Chapapría and C. Viladrich, ``Modelo predictivo agroclimático para detectar cambios usando CRISP-DM,'' Alicia Concytec, 2024. [Online]. Available: \url{https://alicia.concytec.gob.pe/vufind/index.php/Record/UEPU_8c8566d0d4b73311211673ba80756dcb} 

\bibitem{canchano2024} J. C. Canchano, ``Model for demand and price prediction of agricultural products using CRISP-DM,'' in \textit{Proc. IEEE Conf. Data Mining Applications}, 2024, pp. 1--6, doi: 10.1109/ICDMMA.2024.11141314. 

\bibitem{midagri2025} Ministerio de Desarrollo Agrario y Riego, ``Zonificación ecológica económica (ZEE): Suelos y aptitud agrícola,'' 2025. [Online]. Available: \url{https://www.midagri.gob.pe/portal/43-sector-agrario/suelo} 

\bibitem{una2023} Universidad Nacional del Altiplano, ``Propuesta de zonificación agroecológica para el desarrollo sostenible,'' Tesis, 2023. [Online]. Available: \url{https://repositorio.uap.edu.pe/jspui/bitstream/20.500.12990/3250/1/Tesis_propuesta_zonificación_agroecológica_desarrollo_sostenible.pdf} 

\bibitem{chapman2000}
P. Chapman \textit{et al.}, ``CRISP-DM 1.0: Step-by-step data mining guide,'' SPSS Inc., 2000.

\bibitem{midagri2025}
Ministerio de Desarrollo Agrario y Riego, ``Series históricas de producción agrícola -- hoja c-18,'' 2025. [Online]. Available: \url{https://www.midagri.gob.pe/}

\bibitem{nasa2025}
National Aeronautics and Space Administration, ``NASA POWER Data Access Viewer,'' 2025. [Online]. Available: \url{https://power.larc.nasa.gov/}

\bibitem{viladrich}
L. Chapapría and C. Viladrich, ``La relevancia económica de las variables climáticas en la agricultura,'' Centro de Estudios Andaluces. [Online]. Available: \url{https://centrodeestudiosandaluces.es/datos/actividades/Viladrich.pdf}

\bibitem{senamhi2024}
Servicio Nacional de Meteorología e Hidrología del Perú, ``El Niño Costero 2023--2024 fue el más intenso de los últimos 20 años en el oeste de Sudamérica,'' Apr. 25, 2024. [Online]. Available: \url{https://www.gob.pe/institucion/senamhi/noticias/944204-el-nino-costero-2023-2024-fue-el-mas-intenso-de-los-ultimos-20-anos-en-el-oeste-de-sudamerica}

\bibitem{rcr2023}
Redacción RCR Perú, ``Producción agraria en Puno cayó en 70 por ciento desde el 2022 por sequía,'' Jul. 14, 2023. [Online]. Available: \url{https://www.rcrperu.com/produccion-agraria-en-puno-cayo-en-70-por-ciento-desde-el-2022-por-sequia/}

\bibitem{cepiura2023}
Centro de Estudios Piura, ``Nota informativa: Estructura productiva de la región Piura,'' Gobierno Regional Piura, 2023. [Online]. Available: \url{https://www.cepiura.org.pe/wp-content/uploads/2023/04/NOTA-INFORMATIVA-002-2023-2.pdf}

\bibitem{esan2023}
Universidad ESAN, ``La Libertad: ¿Cuál es su aporte a la economía exportadora del Perú?,'' 2023. [Online]. Available: \url{https://www.esan.edu.pe/conexion-esan/la-libertad-cual-es-su-aporte-a-la-economia-exportadora-del-peru}

\bibitem{redesarrolloica2021}
Programa Redesarrollo, ``Dato Regional Ica: Exporta de todo,'' 2021. [Online]. Available: \url{https://www.redesarrollo.pe/multimedia/dato-regional-ica-exporta-de-todo/}

\bibitem{inei1996}
Instituto Nacional de Estadística e Informática, ``Agropecuario: Conociendo Junín,'' 1996. [Online]. Available: \url{http://proyectos.inei.gob.pe/web/biblioineipub/bancopub/est/lib0269/agrope.htm}

\bibitem{dataperu2021}
Programa Redesarrollo, ``San Martín: Economía, salud, educación, hogares, demografía,'' DataPerú, 2021. [Online]. Available: \url{https://data-peru.itp.gob.pe/profile/geo/san-martin}

\bibitem{midagripuno2020}
Ministerio de Desarrollo Agrario y Riego, ``Perfil productivo regional de Puno,'' 2020. [Online]. Available: \url{https://repositorio.midagri.gob.pe/bitstream/20.500.13036/2055/1/Perfil%20productivo%20de%20puno.pdf}

\end{thebibliography}

\end{document}